\PDEIndexMark{idx:hilbert-schmidt-eigenfunction-expansion-introduction}利用这些结果, 可以证明以下特征函数展开定理. 它在求解椭圆边值问题时将非常有用.

\begin{Theorem}{Hilbert-Schmidt 定理}
\label{thm:hilbert-schmidt-spectral}
设 $H$ 为 Hilbert 空间, $A\in\mathcal{L}(H)$ 为紧自伴随算子.

那么存在一个非零实特征值序列 $\{\lambda_i\}_{i=1}^N$, 其中 $N$ 等于算子 $A$ 的秩, 使得 $|\lambda_i|$ 单调不增, 且若 $N=\infty$, 则
\begin{equation}\label{eq:compact-self-adjoint-nonzero-eigenvalue-limit}
\lim_{i\to\infty}\lambda_i=0.
\end{equation}
此外, 若 $A$ 的每个特征值按其重数在序列中重复, 则存在由相应特征函数组成的标准正交集 $\{\phi_i\}_{i=1}^N$; 即
\begin{equation}\label{eq:compact-self-adjoint-eigenfunction-equations}
A\phi_i=\lambda_i\phi_i.
\end{equation}
而且, $\{\phi_i\}_{i=1}^N$ 为 $\mathcal{R}(A)$ 的一组标准正交基; $A$ 可表示为
\begin{equation}\label{eq:hilbert-schmidt-operator-eigenfunction-expansion}
Au=\sum_{i=1}^N\lambda_i(\phi_i,u)\phi_i.
\end{equation}

\end{Theorem}

\begingroup
\postdisplaypenalty=10000
\renewcommand{\endProof}{\par\nobreak\hfill$\blacksquare$\par\vspace{0.1cm}}
\begin{Proof}
注意, 由 Theorem~\ref{thm:symmetric-operator-spectral-properties}, 由于 $A$ 自伴随, $A$ 的特征值都是实数. 由 Fredholm 二择一定理 (Theorem~\ref{thm:fredholm-alternative-compact}), 非零特征值离散, 有界且具有有限重数. 因而可以将它们按绝对值递减的顺序列为序列 $\{\lambda_i\}_{i=1}^N$ (按重数重复), 其中 $N$ 可能为无穷. 由于特征值除零以外不可能有其他聚点, 若 $N$ 为无穷, 则 \eqref{eq:compact-self-adjoint-nonzero-eigenvalue-limit} 必然成立.

现在对每个不同的非零特征值所对应的特征空间选取一组标准正交基, 并将这些基汇集起来 (按照所对应的特征值编号), 组成序列 $\{\phi_i\}_{i=1}^N$. 由 Theorem~\ref{thm:symmetric-operator-spectral-properties}, 整个集合标准正交.

设 $M$ 为 $\{\phi_i\}_{i=1}^N$ 所张成空间的闭包. 断言 $M\supseteq\mathcal{R}(A)$. 注意, 由于 $A$ 自伴随, $M$ 和 $M^\perp$ 均在 $A$ 下不变. 设 $\widehat A$ 为 $A$ 在 $M^\perp$ 上的限制. 由于 $A$ 自伴随且紧, 算子 $\widehat A\in\mathcal{L}(M^\perp)$ 也自伴随且紧. 因而由 Theorem~\ref{thm:fredholm-alternative-compact}, $\widehat A$ 的每个非零谱值都是特征值. 然而, $\widehat A$ 的任何特征值也都是 $A$ 的特征值. 因而 $\widehat A$ 的谱半径为零. 由 Problem~\ref{ex:self-adjoint-spectral-radius-equals-norm}, 这说明 $\widehat A$ 为零算子. 因而 $M^\perp$ 的每个元素都是对应于特征值 $0$ 的特征向量. 所以 $M^\perp=\mathcal{N}(A)$, 且 $\{\phi_i\}_{i=1}^N$ 构成 $\mathcal{R}(A)$ 的一组基.

现在, 由于 $\{\phi_i\}_{i=1}^N$ 构成 $\mathcal{R}(A)$ 的一组基, 有\[\begin{aligned}
A(u)&=\sum_{i=1}^N(\phi_i,A(u))\phi_i\\
&=\sum_{i=1}^N(A(\phi_i),u)\\
&=\sum_{i=1}^N\lambda_i(\phi_i,u)\phi_i.
\end{aligned}\]证明完毕.
\end{Proof}
\endgroup

以下重要推论给出了解非齐次问题的方法.

\begin{Corollary}
\label{cor:compact-self-adjoint-nonhomogeneous-alternative}
设 $A\in\mathcal{L}(H)$ 为紧自伴随算子, $\{\lambda_i\}_{i=1}^N$ 为非零特征值, $\{\phi_i\}_{i=1}^N$ 为前一定理所述的相应特征函数. 对任意 $f\in H$, 令
\begin{equation}\label{eq:hilbert-schmidt-nullspace-data-projection}
f_{\mathcal{N}(A)}:=f-\sum_{i=1}^N(\phi_i,f)\phi_i
\end{equation}
为 $f$ 在 $A$ 的零空间上的投影. 那么对于 $\lambda\ne0$ 的非齐次问题
\begin{equation}\label{eq:compact-self-adjoint-nonhomogeneous-equation}
Au-\lambda u=f,
\end{equation}
有以下二择一结论:\begin{enumerate}\item $\lambda$ 不是 $A$ 的特征值, 此时 \eqref{eq:compact-self-adjoint-nonhomogeneous-equation} 有唯一解
\begin{equation}\label{eq:compact-self-adjoint-noneigenvalue-unique-solution}
u=\sum_{i=1}^\infty\dfrac{(\phi_i,f)}{\lambda_i-\lambda}\phi_i-\dfrac{1}{\lambda}f_{\mathcal{N}(A)};
\end{equation}
或者\item $\lambda$ 是 $A$ 的特征值. 此时, 令 $J$ 为满足 $\lambda_j=\lambda$ 的自然数 $j$ 所组成的有限指标集. 那么 \eqref{eq:compact-self-adjoint-nonhomogeneous-equation} 有解, 当且仅当
\begin{equation}\label{eq:compact-self-adjoint-eigenvalue-solvability-condition}
(\phi_j,f)=0\quad\text{for all }j\in J.
\end{equation}
在这种情形下, \eqref{eq:compact-self-adjoint-nonhomogeneous-equation} 有解族
\begin{equation}\label{eq:compact-self-adjoint-eigenvalue-solution-family}
u=\sum_{j\in J}c_j\phi_j+\sum_{i\in\mathbb{N}\backslash J}\dfrac{(\phi_i,f)}{\lambda_i-\lambda}\phi_i-\dfrac{1}{\lambda}f_{\mathcal{N}(A)},
\end{equation}
其中 $\{c_j\}_{j\in J}$ 为任意常数.\end{enumerate}
\end{Corollary}

\begingroup
\postdisplaypenalty=10000
\renewcommand{\endProof}{\par\nobreak\hfill$\blacksquare$\par\vspace{0.1cm}}
\begin{Proof}
由 Fredholm 二择一定理和 Hilbert-Schmidt 定理, 写出
\begin{equation}\label{eq:compact-self-adjoint-solution-nullspace-decomposition}
u=\sum_{i=1}^N(\phi_i,u)\phi_i+u_{\mathcal{N}(A)},
\end{equation}
展开 \eqref{eq:compact-self-adjoint-nonhomogeneous-equation}, 并比较系数, 即得到证明. 细节留给读者.
\end{Proof}
\endgroup
